\section{Conclusions and discussion}

In this work we have presented a concrete prescription to renormalize non-perturbatively the matrix elements needed
for the computation of quasi-PDFs. The employed scheme is RI$'$, which is then converted to the $\MSb$ scheme and evolved to 2 GeV; this is done
perturbatively to one-loop. We have argued that the renormalization condition properly handles both kinds of divergences present in the
matrix elements: the standard logarithmic divergence and the power divergence specific to non-local operators containing a Wilson line. Furthermore,
we provide the renormalization conditions to eliminate the mixing in the case of the unpolarized quasi-PDF that mixes with the twist-3 scalar operator.

\medskip
We have also demonstrated the implementation of the proposed prescription to the helicity quasi-PDF and presented the corresponding renormalized
matrix elements. This has allowed us to draw conclusions how to make the computation more robust, which is the main outcome of this work.
\begin{itemize}
\item First, an essential ingredient of a computation with controlled systematic uncertainties is the subtraction of one-loop lattice artifacts in the framework
of lattice perturbation theory.
Following the ideas of Ref.~\cite{Alexandrou:2015sea}, we are currently computing the $\mathcal{O}(g^2\,a^\infty)$ artifacts that will be subtracted
from the non-perturbative estimates for the $Z$-factors. In this way, the presence of large cut-off effects in the renormalized functions (especially
for ``parallel'' momenta) will be avoided to a large extent.
\item Second, the conversion factor from the RI$'$ renormalization scheme to the $\MSb$
scheme is likely to have sizable higher order corrections that, among others, are responsible for the unphysical feature of the real part of the renormalized
matrix element becoming negative for large Wilson line lengths. A two-loop computation of this conversion factor is expected to resolve this issue
to a sufficient degree. We have performed numerical experiments that indicate that a natural change of the conversion factor by two-loop contributions,
i.e.\ around 10-20\% (which is approximately $\alpha_s$ at the considered scale), should be enough to suppress the unwanted effect. A perturbative
calculation of the conversion factor to two loops is quite laborious and will be presented separately.
\end{itemize}

\medskip
To summarize, the renormalization program presented in this work together with future improvements that are being pursued,
will provide reliable estimates for the renormalization functions of the Wilson line fermion operators. In this fashion, the obtained
renormalized matrix elements can be used as an input to
calculate the quasi-PDFs and match them to light-front PDFs, which is the main aim of the whole approach.
Apart from the helicity case discussed in this work, we will address the transversity PDF in an analogous manner. For the unpolarized case, one needs
to take into account the mixing with the scalar operator, as explained and numerically demonstrated here.
With this work, we have proposed and discussed a complete renormalization program of the quasi-PDFs, which has been a major
uncertainty prior to this work and constitutes a crucial milestone in connecting lattice QCD results to the light-cone PDFs.
